\section{Avance respecto del plan del Informe 1}
\textit{[Indicar qué actividades se cumplieron, cuáles se retrasaron o cambiaron,
y explicar las dificultades y ajustes de rumbo.]}

\section{Plan de trabajo actualizado}

Habiendo implementado el caso base y evaluado preliminarmente la metodología
de descomposición secuencial, el trabajo futuro se concentrará en mejorar la
coordinación entre las jornadas de los buses y sus requerimientos de recarga.
A partir de las limitaciones identificadas, se proponen las siguientes
cinco etapas:

\begin{enumerate}
    \emergencystretch=2em
    \renewcommand{\labelenumi}{(\roman{enumi})}

    \item \textbf{Mejorar la generación de ventanas de carga.}
    Se evaluará incorporar los requerimientos de batería durante la
    construcción o ajuste de las jornadas, considerando los tiempos
    necesarios para trasladarse al electroterminal y recargar.
    Actualmente, estas oportunidades se determinan mediante una política
    reactiva, lo que puede obligar a dividir jornadas y utilizar buses
    adicionales. Se estudiará una formulación de optimización que permita
    seleccionar ventanas de carga compatibles con los viajes programados,
    buscando reducir la flota necesaria y manteniendo la cobertura de
    todas las expediciones.

    \item \textbf{Ampliar la implementación del MILP de carga.}
    Se extenderán progresivamente las pruebas realizadas en instancias
    reducidas de Los Espinos--Santa Rosa hacia conjuntos mayores de buses
    y otros electroterminales. En cada ejecución se registrarán los tiempos
    de resolución y las brechas de optimalidad, para determinar la escala
    que puede abordarse con los recursos disponibles. Si se presentan
    dificultades computacionales, se evaluarán mejoras de la formulación
    y estrategias de resolución por electroterminal, con el objetivo de
    avanzar hacia la programación de carga de toda la red.

    \item \textbf{Validar y ajustar los supuestos del modelo.}
    Se profundizará el análisis de sensibilidad de las distancias y
    velocidades de los traslados sin pasajeros, los tiempos mínimos entre
    viajes y el radio permitido para el \textit{interlining}. También se
    compararán distintas duraciones de los intervalos del MILP, para
    evaluar cómo la discretización temporal afecta los costos, las
    reservas y el tiempo de cálculo. Estas pruebas permitirán justificar
    mejor las simplificaciones adoptadas y determinar cuáles requieren
    ajustes antes de consolidar la solución final.

    \item \textbf{Evaluar la metodología bajo distintos tipos de día.}
    Una vez validado el funcionamiento para el día laboral, se buscará
    extender el análisis a sábados y domingos, sujeto al tiempo disponible.
    Para ello, se prepararán las expediciones correspondientes a cada tipo
    de día y se ejecutará nuevamente la metodología, considerando sus
    diferencias de horarios y frecuencias. En todos los casos se verificará
    la cobertura de los viajes, la factibilidad de los traslados, los límites
    de batería y la capacidad de carga de los electroterminales.

    \item \textbf{Analizar los resultados y preparar los entregables finales.}
    Se compararán las soluciones obtenidas con el caso base mediante los
    indicadores definidos, distinguiendo el costo total, la flota
    operacional, los buses de reserva y la utilización de los cargadores.
    Se identificarán los beneficios y las limitaciones de cada mejora,
    diferenciando los resultados de instancias reducidas de aquellos
    obtenidos para la red completa. Finalmente, se consolidarán las
    jornadas por bus, los itinerarios de recarga y el desglose de costos,
    junto con las conclusiones del proyecto y la documentación necesaria
    para reproducir los resultados.

\end{enumerate}

\begin{table}[H]
\centering
\caption{Carta Gantt actualizada para los pasos futuros}
\label{tab:gantt-informe2}
\scriptsize
\resizebox{\textwidth}{!}{%
\begin{tabular}{@{}p{0.29\textwidth}p{0.23\textwidth}p{0.18\textwidth}*{6}{c}@{}}
\toprule
\textbf{Actividad / hito}
& \textbf{Depende de}
& \textbf{Responsable(s)}
& \textbf{S1}
& \textbf{S2}
& \textbf{S3}
& \textbf{S4}
& \textbf{S5}
& \textbf{S6} \\
\midrule

1. Mejorar la generación de ventanas de carga
& Metodología y caso base disponibles
& \textit{[Nombre(s)]}
& X & X & X & $\square$ & $\square$ & $\square$ \\

2. Ampliar la implementación del MILP de carga
& Validación inicial de la actividad 1
& \textit{[Nombre(s)]}
& $\square$ & $\square$ & X & X & X & $\square$ \\

3. Validar y ajustar los supuestos del modelo
& Modelo ejecutable; avances de las actividades 1 y 2
& \textit{[Nombre(s)]}
& $\square$ & X & X & X & X & $\square$ \\

4. Evaluar días sábado y domingo
& Modelo laboral validado en las actividades 2 y 3
& \textit{[Nombre(s)]}
& $\square$ & $\square$ & $\square$ & $\square$ & X & $\square$ \\

5. Analizar resultados y preparar los entregables finales
& Resultados progresivos de las actividades 1--4
& \textit{[Nombre(s)]}
& $\square$ & $\square$ & $\square$ & X & X & X \\

Entrega del informe final: 27 de noviembre
& Cierre de la actividad 5
& \textit{[Nombre(s)]}
& $\square$ & $\square$ & $\square$ & $\square$ & $\square$ & X \\

\bottomrule
\end{tabular}%
}
\end{table}